\subsection{1. Kết quả đạt được}

\begin{itemize}[label={--}, itemsep=1pt]
    \item \textbf{Xây dựng thành công kiến trúc đồ thị nhận thức ràng buộc:} Đề tài đã đề xuất và cài đặt hoàn chỉnh hệ thống phát hiện Sybil thông qua việc mô hình hóa dữ liệu Web3 thành đồ thị đa quan hệ với 4 tầng phân biệt (Follow, Interact, Similarity, Co-Owner). Việc áp dụng cơ chế chuẩn hóa trọng số logarit giúp mô hình giải quyết hiệu quả bài toán nhiễu và bùng nổ thông điệp tại các nút siêu đặc trên đồ thị mạng xã hội phi tập trung.
    
    \item \textbf{Tối ưu hóa kiến trúc và đột phá về kết quả phân lớp:} Kết quả thực nghiệm chứng minh kiến trúc phân lớp tách rời (\textit{Decoupled GNN-ML}) mang lại hiệu năng vượt trội. Tổ hợp giữa GATv2 (đóng vai trò trích xuất đặc trưng) và Random Forest (bộ phân lớp) đã đạt chỉ số Macro F1-Score lên tới 0.9303 trên tập dữ liệu khó nhất (phạm vi Quý). Đặc biệt, mô hình đạt tỷ lệ nhận diện nhầm 0\% đối với nhóm người dùng thật (\textit{Benign}), bảo vệ tuyệt đối trải nghiệm của người dùng chân chính trong khi vẫn tóm gọn hiệu quả các mạng lưới Sybil lớn.
    
    \item \textbf{Giải quyết bài toán thiếu hụt nhãn:} Đề tài đã thiết kế thành công pipeline 3 giai đoạn, tự động hóa hoàn toàn quy trình gán nhãn dữ liệu bằng học không giám sát với GAE kết hợp hệ luật suy diễn.
\end{itemize}